% Generated draft: requires mathematical and source review.
\documentclass[11pt]{article}
\usepackage[margin=25mm]{geometry}
\usepackage{fontspec}
\setmainfont{Latin Modern Roman}
\usepackage{amsmath,amssymb,mathtools,unicode-math}
\setmathfont{Latin Modern Math}
\usepackage{enumitem,needspace,xurl,hyperref}
\setlength{\parindent}{0pt}
\setlength{\parskip}{5pt}
\setlength{\emergencystretch}{3em}
\newcommand{\N}{\mathbb{N}}
\newcommand{\Z}{\mathbb{Z}}
\newcommand{\Q}{\mathbb{Q}}
\newcommand{\R}{\mathbb{R}}
\newcommand{\C}{\mathbb{C}}
\newcommand{\F}{\mathbb{F}}
\newcommand{\A}{\mathbb{A}}
\newcommand{\PP}{\mathbb{P}}
\newcommand{\E}{\mathbb{E}}
\newcommand{\eps}{\varepsilon}
\newcommand{\dd}{\,\mathrm{d}}
\begin{document}
\section*{UnsolvedMath ID: 30001428}
\textbf{Source problem code:} catalogue entry
\subsection{Definitions and mathematical statement}
The supplied record concerns an infinite random graph \(G\) described as hyperbolic, with finite random horizontal levels joined by a perpendicular spanning tree. It does not specify the probability law, the indexing and sizes of the levels, the horizontal edge distribution, the attachment of the spanning tree, or the meaning of perpendicularity. Consequently this description does not determine a random graph model.

Write \(B(G)\) provisionally for the boundary intended in the source, equipped with its intended topology. Neither the boundary construction nor its topology is specified in the supplied passage; in particular, identifying it with the Gromov boundary would require confirmation. A topological space \(X\) is path-connected if, for every \(x,y\in X\), there exists a continuous map \(\gamma\colon[0,1]\to X\) satisfying \(\gamma(0)=x\) and \(\gamma(1)=y\). A finite graph is connected if every pair of its vertices can be joined by a finite graph path.

\textbf{Question as supplied, with unresolved qualifications.} For the intended class of infinite random hyperbolic graphs built from finite random horizontal levels joined by a perpendicular spanning tree, determine the structure of the random boundary \(B(G)\). Is there a threshold for path-connectedness of \(B(G)\), and how is it related to the threshold for connectivity of the horizontal levels?

This is an informal question rather than a fully quantified conjecture. The supplied record does not identify a varying parameter or limiting regime, define either threshold, specify the probability quantifiers, or state the proposed relationship between the thresholds. It also does not say whether connectivity of horizontal levels means connectivity of an individual level, every level, all sufficiently distant levels, or an asymptotic probability of connectivity. These alternatives remain unresolved; no particular interpretation or equality of thresholds is asserted here.
\subsection{Short English statement}
The question asks what the boundary of a random hyperbolic graph built from finite horizontal levels and a spanning tree looks like, and whether its path-connectedness has a threshold related to connectivity within the levels. The graph model, boundary, probability quantifiers, and meaning of threshold must be specified before a precise conjecture can be stated.
\subsection{Sources}
\begin{itemize}
\item Supplied local UnsolvedMath registry record, permanent ID \(30001428\), title \emph{Boundaries of Random Hyperbolic Level Graphs}, source problem code \texttt{catalogue entry}, in the \emph{Oberwolfach Reports Open Problems} collection. The supplied proposed year is 2010; no individual proposer is named.
\item Supplied source citation: \emph{Graph Theory} (2010), \emph{Oberwolfach Reports}, DOI \texttt{10.4171/owr/2010/11}, pp.~521--580. Source URL: \url{https://doi.org/10.4171/owr/2010/11}. The background identifies Oberwolfach Reports as records of mathematical workshops and their open-problem discussions. No page locating this particular problem within the report is supplied.
\item Supplied literature source: \emph{Oberwolfach Report 11/2010}, \url{https://doi.org/10.4171/owr/2010/11}.
\item Literature assessment dated \texttt{2026-08-21}: the supplied curation labels the problem open, describes the extracted passage as an informal random-boundary problem, and reports that the cited report gives no resolution and that no later primary resolution was located. This is supplied dated curation, not an independently verified literature assessment; time-sensitive status and cited claims require independent verification.
\item The registry labels the statement \texttt{exact} and reports that its mathematical content was checked against the official report, with notation, context, and question form normalized without changing the source-level claim. These are supplied verification claims; the report has not been independently consulted for this draft.
\end{itemize}
\end{document}
